\section{Method details}
\label{sec:app_method}

This appendix holds the detail of Section~\ref{sec:method} that the main text summarises.

\begin{figure}[htbp]
\centering
\begin{tikzpicture}[>=Stealth, font=\small,
    blk/.style={draw, thick, rounded corners=2pt, align=center, minimum height=0.85cm, font=\small},
    ours/.style={blk, fill=blue!6, draw=blue!50!black},
    his/.style={blk, fill=orange!10, draw=orange!60!black},
    fc/.style={blk, fill=gray!12, draw=gray!60},
    lab/.style={font=\scriptsize, align=center}]
  \node[blk]  (op)   at (5.6,1.7)   {Operator\\ \scriptsize disarm button};
  \node[ours] (fm)   at (5.6,0)     {Fleet manager\\ \scriptsize phases, faults, LAND};
  \node[ours] (plan) at (0,0)       {Load--rod planner\\ \scriptsize 10\,Hz; owns the ring};
  \node[ours] (trk)  at (0,-2.5)    {Tracker $i$ (carrier)\\ \scriptsize 50\,Hz; owns UAV $i$};
  \node[ours] (trkn) at (5.6,-2.5)  {Tracker (newcomer)\\ \scriptsize from the attach, or take-back};
  \node[his]  (tej)  at (11.3,-2.5) {Collaborator's controller\\ \scriptsize pickup, transit, drop,\\ \scriptsize descent to contact, lift};
  \node[ours] (mux)  at (5.6,-4.6)  {Mux\\ \scriptsize one owner, latched};
  \node[fc]   (bf)   at (0,-4.6)    {Flight controller $i$\\ \scriptsize Betaflight rate loop, magnet};
  \node[fc]   (bfn)  at (5.6,-6.3)  {Flight controller (newcomer)\\ \scriptsize Betaflight rate loop, magnet};
  \draw[->, thick] (op) -- node[lab, right] {commands} (fm);
  \draw[->, thick] (fm) -- node[lab, above] {phase} (plan);
  \draw[->, thick] (plan) -- node[lab, left] {reference,\\ rod pull} (trk);
  \draw[->, thick] (plan.south east) -- node[lab, sloped, above, pos=0.55] {once folded in} (trkn.north west);
  \draw[->, thick] (trk) -- node[lab, left] {throttle, rates,\\ arm, magnet} (bf);
  \draw[->, thick] (trkn) -- (mux);
  \draw[->, thick] (tej.south) -- ++(0,-1.25) -- node[lab, above] {until the attach} (mux.east);
  \draw[->, thick] (mux) -- node[lab, right] {one command stream} (bfn);
  \draw[->, thick, dashed] (trk) -- node[lab, above] {fault} (trkn);
  \draw[->, thick, dashed] (fm) -- node[lab, right] {LAND, drop, disarm} (trkn);
  \draw[->, thick, dashed] (fm.north west) -- ++(0,0.5) -| node[lab, pos=0.25, above] {arm, disarm, LAND} ([xshift=-0.5cm]plan.north west) |- (trk.west);
  \draw[<->, thick, dotted] (fm.east) -| node[lab, pos=0.25, above] {attach confirmed, give-up} (tej.north);
\end{tikzpicture}
\caption[The cascade and its owners]{The cascade and its owners. Blue: this work; orange: the collaborator's
controller; grey: the unmodified flight controllers. The planner owns the ring and the carriers' references; each
tracker owns one UAV's radio link; the newcomer's link is owned by the collaborator's controller until the attach
is confirmed and by this work's tracker from that tick, through a multiplexer that latches. Dashed: faults and the
fleet manager's commands; dotted: the confirmation and the give-up exchanged with the collaborator's stack. The
operator's disarm overrides everything.}
\label{fig:owners}
\end{figure}

\subsection{The tracker's problem and the model integral in detail}
\label{sec:app_integral}

The tracker solves, every 20\,ms,
\begin{equation}
\begin{aligned}
    \min \quad & \sum_{k=0}^{N_h-1} \big\| \mathbf{h}(\mathbf{x}_k, \mathbf{u}_k) - \mathbf{h}_k^{\mathrm{ref}} \big\|_{Q}^{2} + \big\| \mathbf{h}_{N_h}(\mathbf{x}_{N_h}) - \mathbf{h}_{N_h}^{\mathrm{ref}} \big\|_{Q_{N_h}}^{2} \\
    \text{s.t.} \quad & \mathbf{x}_{k+1} = \mathbf{F}_i(\mathbf{x}_k, \mathbf{u}_k; \mathbf{a}_{c,i,k}), \qquad \mathbf{x}_0 = \hat{\mathbf{x}}_i(t), \\
    & 0.05 \le u_{T,k} \le u_{T,\max}, \quad |c_{\phi,k}|, |c_{\theta,k}|, |c_{\psi,k}| \le 1, \quad |\mathbf{u}_k| \le \overline{\mathbf{u}},
\end{aligned}
\label{eq:trk_ocp}
\end{equation}
with $\mathbf{h} = [\mathbf{p}, \mathbf{v}, \boldsymbol{\omega}, \mathbf{c}, \mathbf{u}, \mathbf{e}_R]$, the
attitude reference being the direction of the planned thrust \eqref{eq:thrust_req}. The stick states
are normalised to $[-1, 1]$, and $\overline{\mathbf{u}}$ limits how fast each may change (throttle
0.5\,s$^{-1}$). Only the first command of each solution is sent.




\paragraph{Properties.} At an equilibrium with the gate open ($G = I$) and no component of $\mathbf{b}$
saturated, $\mathbf{b}_{k+1} = \mathbf{b}_k$ gives $\mathbf{e} = \mathbf{0}$, whatever the constant cause:
by the internal model principle \cite{francis1976internal}, the integrator is the model of a constant.
This assumes stability; it does not give it. Approximating the cascade's response to $\mathbf{d}$ by a
first-order lag with time constant $\tau$, each axis becomes
\begin{equation}
    \tau s^2 + s + k_i = 0, \qquad \zeta = \frac{1}{2\sqrt{k_i \tau}},
\end{equation}
so a faster integral costs damping, and delays through the planner and the radio link reduce the margin
further. Let an unmodelled constant force
$\mathbf{F}$ act on the payload and let the cascade, of actual stiffness $K_c$, be quasi-static (fast
compared with the integral). The plan cancels the force it believes in, so the payload settles at
$\mathbf{e} = -K_c^{-1}(\mathbf{F} - \mathbf{d}) = -K_c^{-1}(\mathbf{F} + K\mathbf{b})$, and in continuous time
\begin{equation}
    \dot{\mathbf{b}} = k_i\,\mathbf{e} = -k_i K_c^{-1}\big(K\mathbf{b} + \mathbf{F}\big)
    \quad\Rightarrow\quad \mathbf{b} \to -K^{-1}\mathbf{F}, \qquad \mathbf{d} \to \mathbf{F},
\end{equation}
a first-order convergence at the rate $k_i K K_c^{-1}$ on each axis, with time constant $1/k_i = 2.5$\,s when
$K = K_c$. The believed force converges to the applied one whatever $K$ is: $K$ sets only the speed, and
too large a $K$ costs damping as too large a $k_i$ does. When the force is removed, $\mathbf{b}$ still holds
$-K^{-1}\mathbf{F}$, so the payload overshoots by $K_c^{-1}\mathbf{F}$, the displacement the force caused at
onset, and the overshoot decays at the same rate. This analysis assumes that the force reaches the payload's
position error. With the measured rod pull (Section~\ref{sec:tracker}) a constant vertical force is taken up at
each UAV within about a second, before the integral can act, so in the twin the integral is left with the
horizontal axes (Section~\ref{sec:res_force}).

\paragraph{Reference-shift form.} An earlier version added the height integral to the payload's height
target instead, $z_L^{\mathrm{ref}} \leftarrow z_L^{\mathrm{ref}} + b_z$. Both forms reach the same steady
state; in offset-free MPC they are output and input disturbance models, which differ in the transient, not
in the steady state \cite{pannocchia2007combined}. The model form keeps the target fixed, so its plan
predicts a payload held at the target, which is what happens; the reference-shift form plans a move towards
a shifted target that never comes. The model form also extends to the horizontal axes, and it is the
default everywhere, a resizing fleet included (the integral's force is kept through the resize). The
reference-shift form, on the height only, flies where the model form cannot, and flew on the rig on 1 and
8 October ($k_i = 0.4$\,s$^{-1}$); the 7 October carry flights flew the model form.
% Model integral through a resize since 8 Oct (896aff2; bringup real_mode.planner_options).
% STATUS 5 Oct 2026: int_mode auto (model where allowed, else reference): reference with kt_trim on,
% a resizing fleet, z_ki 0, offset_free on or dist_est ring/full (bringup real_mode.planner_options).

\paragraph{Bound and limits.} $\bar{b} = 0.3$\,m bounds the force the integral can supply to $K\bar{b}$
(2.1\,N vertically), enough for an unmodelled 0.1\,kg object. A larger error is a fault, not a disturbance,
and a warning at 70\,\% of the bound reports it rather than covering it silently; the bound and the warning are
also what stop an integral that winds up without effect, as the vertical integral does on a wrong attachment
radius in the twin (Section~\ref{sec:sim_robustness}). A force on the payload cannot represent a moment: one UAV's
thrust-map error is a moment on the ring, which that UAV's own tracker takes up through its measured pull. A miss
larger than the gate is never integrated. The integral is also slow ($1/k_i = 2.5$\,s): it removes constant and
slowly varying forces, not impacts or gusts. Five other ways to remove a steady error were built and flown on the
earlier stack (a force balance \cite{muske2002disturbance}, an offset-free observer
\cite{pannocchia2003disturbance}, the measured cable force, a per-UAV thrust-gain trim, and the rods' mass in the
planner); Appendix~\ref{sec:app_methods} records them, and the measured cable force became the trackers' measured
pull.

\subsection{Reconfiguration: the ramp, the seed, the fall-backs and the constants}
\label{sec:app_reconfig}

\paragraph{The ramp and the seed.} Where the planner is told of a change it can shape the step. For a ramped
release the leaver's planned tension is run down and the survivors' run up over a short ramp before the magnet is
released, in place of a step at the resize tick; the first version of the ramp, which ran the leaver's tension
down but let the survivors' step, is the negative that defined it. For a join the newcomer is seeded at 0.1\,N and
every tension is blended to the balanced solution for the new geometry, as a dock blends thrust across its
controller switch \cite{sugihara2023trady}: on the earlier plant a newcomer started at a full share pulled the ring
with a lever about the incumbents' rod apex and tilted it at the weld (Appendix~\ref{sec:app_proximity_weld}).



\paragraph{The freed UAV.} After a release the freed UAV steps 0.5\,m out radially at its height, clear of the ring
and the survivors, and lands by itself while the survivors carry on: its height reference is pressed below its own
resting height until it is there, and it is then disarmed alone, without waiting for the fleet
(Section~\ref{sec:res_freed}). Touchdown for the fleet is judged on the survivors only, so a freed UAV still
descending elsewhere does not hold their references down, and the fleet is declared landed only when it is down
too; the rod feedforward is switched off once the ring rests, because its sideways component pushes a grounded
survivor across the floor.

\paragraph{The take-back.} A newcomer whose contact is not confirmed within its timeout, or whose controller gives
up, is taken back: this work's tracker latches onto its measured position and velocity, ramps its reference at
0.25\,m/s$^2$ so that the latch makes no step, lifts it clear of the rim and the carriers, moves it out and lands it
by itself. On the circle the hold moves with the ring in the ring's planned frame, so that an orbiting formation
cannot sweep a carrier into the newcomer's rod; where the newcomer is low and far out, it lands where it stops. On
the rig the take-back first rises only 5\,cm above the tether's slack and waits for the tip to leave the plate, so
that a magnet that has caught after all is not torn off; if the plate rises instead, the magnet is switched off
again and the fleet lands. The tilt cut, the drop rule and LAND are the exits of Section~\ref{sec:safety}.

\paragraph{The rejoin mission.} The full mission is assembled with the collaborator's mission planner. Four UAVs
carry the ring in hover or round the circle; the one on plate 9 is released and flies the collaborator's errand
(picking up an object with its magnet and dropping it into the ring) and comes back to its plate; it is admitted
under the contract above, folded in, and the fleet lands. On the rig the collaborator's controller does every
attach down to contact, since the magnet catches only on the plate; this work's own approach, which steered the
measured magnet tip and welded on proximity, flies in the twin only (Appendix~\ref{sec:app_proximity_weld}).


\begin{table}[htbp]
\centering
\small
\caption[Reconfiguration constants and their reasons]{The constants of the reconfiguration and the reason for each
value: a noise floor times a factor, or a physical bound. [Bracketed reasons wait for the offline sweeps over the rig
recordings.]}
\label{tab:constants}
\begin{tabularx}{\textwidth}{>{\raggedright\arraybackslash}p{3.6cm}>{\raggedright\arraybackslash}p{2.4cm}>{\raggedright\arraybackslash}X}
\toprule
Constant & Value & Reason \\
\midrule
Rod stretch for an unannounced leave & 0.06\,m, two ticks & A rod on its plate cannot read longer than itself; the
threshold sits above the scatter of a rod length measured from motion capture, and two ticks reject a single bad
pose. [False alarms over the rig recordings and delay against threshold from the rig releases.] \\
Contact & catch radius 24\,mm (rig); face $-15$ to $+12$\,mm; under 0.03\,m/s for 0.3\,s & The magnet catches only on
the plate: the tip must be over the plate's face and at it, and still for 36 motion-capture samples, which excludes a
pass over the plate. \\
Confirmation, 1.0\,s after the lift & tip within 6\,mm of its contact point; plate range 5\,mm or rise 3\,mm; taut
within 3\,mm & A caught tip moves with the plate and a missed one is left behind by it; the free hang is measured on
each flight, so the taut check reads this flight's tether. \\
Confirmation fails & 45\,mm for 0.12\,s; any data gap over 0.1\,s restarts & A tip that leaves the plate has not
caught; no hold is judged across missing poses. \\
Newcomer's seed tension & 0.1\,N & Small against the newcomer's balanced share (1.6\,N on plate 9 of the carry
layout), so the fold-in's blend and not the weld sets the step. \\
Step-out after a release & 0.5\,m radially & Clears the rim (0.255\,m) and the survivors' rods from the leaver's
plate. \\
Take-back ramp & 0.25\,m/s$^2$ & A quasi-static tilt of 1.5$^\circ$ on the newcomer, under the 2$^\circ$ bar on its
attitude step, and it stops 0.34\,m/s in 1.4\,s. \\
Take-back clearances & 0.15\,m over the carriers; 0.3\,m under the ring; 0.9\,m from the ring's centre & The
newcomer's body and rod stay a rotor disc clear of the carriers; a low newcomer lands where it is rather than
rising through the formation. \\
Ring tilt that drops the ring & 30$^\circ$ & Beyond the planner's near-level model and about where a rod's pull turns
across its plate; where the controller gives up, not a value to design to (Section~\ref{sec:safety}). \\
Fewest carriers & 3 & Two rim points make the ring a pendulum about the line between them. \\
\bottomrule
\end{tabularx}
\end{table}

\subsection{Metrics in full}
\label{sec:app_metrics}

\begin{description}
\item[Release metrics] For a UAV leaving the team, frozen on 11 October before the first flight they
    score; one rule for the twin and the rig. The ring's tilt is kept as a vector: with $\mathbf{n}$ the ring's
    body $z$ axis in the world frame, the tilt angle is $\arctan\!\big(\sqrt{n_x^2+n_y^2}\,/\,n_z\big)$ and
    the tilt vector $\boldsymbol{\theta}$ is that angle along the horizontal direction of $\mathbf{n}$, so that
    $|\boldsymbol{\theta}|$ is the scalar tilt above. The event time $t_e$ is the DETACH command for an announced
    release, the magnet's OFF for an unannounced one (the runner's RELEASE in the twin, the first magnet OFF in
    the rig's bag), and the first planner tick at which the leaver's planned tension falls for a ramped release;
    LAND is the LAND command. Over the pre-window $[t_e-5\,\mathrm{s},\,t_e)$ the mean
    $\bar{\boldsymbol{\theta}}_{\mathrm{pre}}$ and the spread
    $\sigma_{\mathrm{pre}}=\sqrt{\overline{|\boldsymbol{\theta}-\bar{\boldsymbol{\theta}}_{\mathrm{pre}}|^2}}$
    are taken. A row whose ring reference was not still for the 5\,s before the event (its height within
    5\,mm of its value at $t_e$ and its speed within 0.02\,m/s of its median over the last 2\,s) is flagged and
    still reported. \textbf{Tilt rise} $\Delta\theta$, the primary metric, is the maximum of
    $|\boldsymbol{\theta}(t)-\bar{\boldsymbol{\theta}}_{\mathrm{pre}}|$ over $[t_e,\,\mathrm{LAND}]$; when two
    runs are compared both are cut at the shorter hold, each 8 October pair at the rig flight's event-to-LAND,
    and the ledger's bands and the flights they score at $t_e+15$\,s. Beside it: the absolute peak,
    $\max|\boldsymbol{\theta}|$ over the same window, and the scalar rise, the absolute peak minus the mean of
    $|\boldsymbol{\theta}|$ over the pre-window. The 30$^\circ$ drop rule and the join bars stay on the
    absolute tilt. \textbf{Time back in band:} with $\bar{\boldsymbol{\theta}}_{\mathrm{post}}$ the mean over
    the last 5\,s before LAND, scored only when LAND comes at least 10\,s after $t_e$, and the half-width
    $b=\max(2\sigma_{\mathrm{pre}},\,0.5^\circ)$, it is the earliest $t$ at or after the time of $\Delta\theta$
    such that $|\boldsymbol{\theta}-\bar{\boldsymbol{\theta}}_{\mathrm{post}}|\le b$ for every sample in
    $[t,\,t+1\,\mathrm{s}]$, reported as $t-t_e$, or ``not back''. The dip (the largest drop of the ring's
    height below its pre-event mean), the detection time (from the magnet's OFF to the planner's resize), the
    steady offset, the landing lean, the clearances and the survivors' shares keep their scorers and are not
    part of the freeze.
    \item[Join metrics] For a newcomer admitted to the ring: the time from the last contact to the confirmation; the ring's
    tilt rise, on the release metric above with the contact as the event, from the contact to 8\,s after the
    confirmation; the ring dip over the same window; the hand-over step, the change in the newcomer's applied
    throttle, roll, pitch and yaw-rate commands between the 0.1\,s before the switch of authority and the 0.3\,s
    after it, as a fraction of its hover throttle and in degrees; the newcomer's tilt at the confirmation; and the
    clearances of the newcomer's body and rod to every carrier body and rod through its approach and, when it is
    taken back, through its take-back. For the full mission, also the tilt at the release, the mean tilt over the
    circle and the time from TAKEOFF to every UAV landed.
    \item[Magnet margin] For a release, the survivors' peak measured rod pull and the rate at which it rises,
    against the magnet's hold at that rod angle (Section~\ref{sec:hardware}); the one metric that says how close a
    transient came to letting a second magnet go.
    \item[Switch validation] For every resize, the planner's solver status and residual at the first live solve
    after the switch, the distance of that solve from an offline solve of the same logged problem run to
    convergence (not computed so far: the planner keeps a problem only when its solve fails), the solve time against the 60\,ms budget on the rig, and the step in each survivor's commanded
    throttle across the switch.
    \item[Solver health] The number of planner solves that failed or did not converge.
\end{description}

Every arm of a lever table and every ledger row is labelled a \emph{replication} when both arms of a twin pair, or the
rig flight of a ledger row, were flown and registered before the rule was signed on 11 October, and a \emph{new
prediction} otherwise; the label is set at the signature and never changes.

\paragraph{Why these bars.} Three physical limits say what a transient may do to the ring, and the bars of the
results sections sit inside them. The first is the magnets. A plate holds its magnet straight off the face by
about 9--10\,N and by less under shear, so a tilt that turns a rod's pull across its plate can let a magnet go:
every ring lost on the rig began with a magnet letting go with no command, two with the rods planned at
45$^\circ$ and two since the rods have been planned at 60$^\circ$ to keep each pull closer to its plate's axis
(Section~\ref{sec:res_leave}). [The hold against the angle of pull, still to be measured with a spring scale,
gives this limit its number.] The second is the fleet's own rule:
a ring tilt over 30$^\circ$, or a loss that leaves fewer than three carriers, drops the ring and lands every UAV
(Section~\ref{sec:safety}), so
30$^\circ$ is where the controller gives up, not a value to design to. The third is the model: the planner's
load--rod model and each tracker's rod pull assume a near-level ring, so a transient that stays within a few
degrees stays inside the regime the plan was computed for, and a transient that leaves it is being flown on a
plan that no longer describes the ring. The bars are therefore stated as a rise above the ring's own motion
before the event, small against the first two limits and inside the third, and every run is reported with its
measured rise so that a run that misses a bar is still read against the limits rather than against a number.
